\scanpage{8}
\mspage{4}
une application
\[
(22)\qquad\widetilde X'(i)\xrightarrow{\Phi_i}\operatorname{Hom}(\mathbf T(\overline{\mathbf Q}),\pi_1(X'))
\]
\margin{\unsure{ici il est supposé} $\dim X=1$ -- autrement il
faut \unsure{voir|revoir} la définition de $\widetilde X'(I)$\ldots\
\unsure{Autre façon de dire : il faut que} $\widetilde X'$ \unsure{soit considéré}
\unsure{pour donner des} \illegible{1} \unsure{henséliens ou locaux au} dessus d'un \unsure{pt de} $X'(I)$\ldots}
compatible avec les opérations de $\pi_1(X')$
de part et d'autre, et qui précise
la donnée de $\varphi_i$ (17). Mieux encore,
le système $(\Phi_i)_{i\in I}=\Phi$ :
\[
(23)\qquad\widetilde X'(I)\xrightarrow{\Phi}\operatorname{Hom}(\mathbf T(\overline{\mathbf Q}),\pi_1(X'))
\]
est compatible avec les $\mathcal E_{X'}\deleted{\unsure{\mathcal E}}$ de part et
d'autre, en se rappelant, pour l'opération
de $\mathcal E_{X'}$ sur le mb de droite, de tenir
compte aussi de l'opération de $\mathcal E_{X'}$
sur $\mathbf T(\overline{\mathbf Q})$ (via $\Gamma$).

\textcircled{5} Soit enfin, par dessus le marché, une
partie $J\subset X'(\overline{\mathbf Q})$. Soit $\widetilde X'(J)$ son
image inverse dans $\widetilde X'$. Soit $\mathcal G_{X'}(I,J)$
ou simplement $\mathcal G'$ le sous-groupe de $\mathcal G_{X'}(I)$
qui invarie $J$. \deleted{\unsure{On voit les structures}} \inserted{Donc $\mathcal G'$}
opère sur $J$, sur $\widetilde X'(J)$ et sur
\[
(24)\qquad\widetilde X'(J)\longrightarrow J
\]
\mspage{7}
